\chapter{Resultados y evaluación}
\label{cap:resultados}

% Todas las cifras de este capítulo proceden de tesis/cifras.tex, que genera
% scripts/generar_cifras.py a partir de los archivos de outputs/.

\section{Resultados, análisis e interpretación}

\subsection{Conjunto analizado}

El procedimiento se ejecutó sobre las \cNEscenas{} imágenes de la cámara de navegación de
\textit{Curiosity} que disponen de máscara de entrenamiento. La
Figura~\ref{fig:paso-a-paso} ilustra el flujo completo sobre una escena representativa:
imagen original, máscara de AI4Mars, binaria depurada, transformada de distancia con las
semillas detectadas y resultado del conteo.

\begin{figure}[H]
  \centering
  \includegraphics[width=\textwidth]{Figura_11_procedimiento_paso_a_paso.png}
  \caption[Procedimiento paso a paso]{Procedimiento paso a paso sobre una escena con
    varias rocas: imagen original, máscara de AI4Mars, binaria depurada, transformada de
    distancia con las semillas y resultado del conteo.}
  \label{fig:paso-a-paso}
\end{figure}

Cada imagen recibió una bandera de calidad que resume su aptitud para cada indicador
(Tabla~\ref{tab:banderas}, Figura~\ref{fig:banderas}).

\begin{table}[H]
  \centering
  \caption{Composición del conjunto según la bandera de calidad.}
  \label{tab:banderas}
  \begin{tabular}{llrr}
    \toprule
    \textbf{Bandera} & \textbf{Significado} & \textbf{Imágenes} & \textbf{\%} \\
    \midrule
    \texttt{ok}          & Contiene roca grande; apta para ambos indicadores & \cFlagOk     & \cFlagOkPct \\
    \texttt{no\_bigrock} & Contiene roca, pero sin roca grande que contar    & \cFlagNoBig  & \cFlagNoBigPct \\
    \texttt{no\_rock}    & Etiquetada, sin roca                             & \cFlagNoRock & \cFlagNoRockPct \\
    \texttt{mostly\_null}& Más del \SI{95}{\percent} sin etiqueta           & \cFlagNull   & \cFlagNullPct \\
    \texttt{empty}       & Sin ningún píxel etiquetado                      & \cFlagEmpty  & \cFlagEmptyPct \\
    \bottomrule
  \end{tabular}
\end{table}

\begin{figure}[H]
  \centering
  \includegraphics[width=0.82\textwidth]{Figura_12_banderas_calidad.png}
  \caption{Composición del conjunto analizado según la bandera de calidad.}
  \label{fig:banderas}
\end{figure}

Las banderas fijan las dos poblaciones de análisis, que se mantienen en todo el documento:

\begin{itemize}
  \item \textbf{Población de E1}: las \cNCobCalc{} escenas con al menos un píxel etiquetado,
    es decir, todas salvo las \cFlagEmpty{} de bandera \texttt{empty}.
  \item \textbf{Población de E2}: las \cFlagOk{} escenas de bandera \texttt{ok}. Quedan fuera
    las \cNNullBig{} escenas de bandera \texttt{mostly\_null} que contienen algún píxel de roca
    grande —y que aportarían \cRocasNullBig{} rocas—, porque cuando más del
    \SI{95}{\percent} de la escena carece de etiqueta la región de roca grande es un fragmento
    aislado cuya forma no puede interpretarse. La subsección~\ref{sec:diagnostico} muestra que
    los casos de anotación anómala identificados pertenecen precisamente a ese grupo.
\end{itemize}

La fracción de escena efectivamente etiquetada tiene una mediana de \cFracValidMed: algo más
de la mitad de cada imagen recibió etiqueta, y el resto corresponde al cuerpo del rover, al
cielo y a las distancias que el propio conjunto excluye por diseño.

\subsection{Cobertura de roca visible (E1)}

La cobertura pudo calcularse en \cNCobCalc{} de las \cNEscenas{} escenas (\cPctCobCalc);
\cNCobPos{} escenas (\cPctCobPos{} del total) presentaron cobertura de roca mayor que cero.
Entre estas últimas, la mediana de cobertura sobre píxeles etiquetados, según la
ecuación~\eqref{eq:cobertura}, fue del \cCobMedVal{} y la media del \cCobMediaVal; medida
sobre la imagen completa, la mediana es del \cCobMedTot{} y la media del \cCobMediaTot.

La distribución (Figura~\ref{fig:dist-cobertura}) es marcadamente \textbf{bimodal}: se
acumula en los extremos, con \cNCobCien{} escenas —el \cPctCobCien{} de las que contienen
roca— en las que la totalidad del área etiquetada es roca, y un segundo grupo de cobertura
nula. En total, \cNCobMitad{} escenas superan el \SI{50}{\percent} de cobertura.

\begin{figure}[H]
  \centering
  \includegraphics[width=\textwidth]{Figura_13_distribucion_cobertura.png}
  \caption[Distribución de la cobertura]{Distribución de la cobertura de roca visible en
    sus dos versiones: sobre píxeles etiquetados y sobre la imagen completa.}
  \label{fig:dist-cobertura}
\end{figure}

\subsubsection{Control del denominador}

La cobertura divide por los píxeles etiquetados. Si las escenas poco etiquetadas tendieran a
conservar sobre todo la roca, sus coberturas serían altas por construcción y no por el
terreno. Para examinarlo se relacionó la cobertura con la fracción de escena etiquetada
(Figura~\ref{fig:cob-vs-frac}) mediante las medidas de dependencia de la
Sección~\ref{sec:evaluacion}.

\begin{figure}[H]
  \centering
  \includegraphics[width=0.85\textwidth]{Figura_14_cobertura_vs_fraccion_etiquetada.png}
  \caption{Cobertura frente a fracción de escena etiquetada.}
  \label{fig:cob-vs-frac}
\end{figure}

No se observó una asociación lineal global entre la cobertura y la fracción etiquetada
($r = \cDepFCPearson$, IC~95\,\% \cDepFCPearsonIC), por lo que no se encontró evidencia de
que las coberturas altas estén explicadas simplemente por una disminución lineal de la
fracción válida. Con \cDepFCN{} escenas, el intervalo excluye el cero: la asociación es
detectable, pero de una magnitud que explica menos de una milésima parte de la varianza. Las
medidas no lineales registran también una dependencia débil —correlación de distancia de
\cDepFCDcor, $p = \cDepFCDcorP$— que la Tabla~\ref{tab:dependencia} recoge junto al resto.
Este análisis no descarta otras formas de dependencia ni los sesgos propios del proceso de
anotación.

La descripción por tramos de fracción etiquetada (Tabla~\ref{tab:cob-por-frac}) aclara la
forma de esa dependencia. Si el artefacto del denominador operase, la cobertura mediana
debería ser máxima en las escenas menos etiquetadas y decrecer con la fracción. No ocurre:
crece hasta el tercer tramo y cae en el último, y las escenas con menos de un cuarto de
superficie etiquetada tienen una cobertura mediana inferior a la de los dos tramos
centrales.

\begin{table}[H]
  \centering
  \caption[Cobertura por tramos de fracción etiquetada]{Cobertura mediana por tramos de
    fracción de escena etiquetada, en sus dos versiones.}
  \label{tab:cob-por-frac}
  \begin{tabular}{lrcc}
    \toprule
     & & \multicolumn{2}{c}{\textbf{Cobertura mediana}} \\
    \cmidrule(lr){3-4}
    \textbf{Fracción etiquetada} & \textbf{Escenas} & \textbf{sobre lo etiquetado}
      & \textbf{sobre la imagen} \\
    \midrule
    hasta 0,25     & \cFTNA & \cFTMedA & \cFTMedTotA \\
    0,25 -- 0,50   & \cFTNB & \cFTMedB & \cFTMedTotB \\
    0,50 -- 0,75   & \cFTNC & \cFTMedC & \cFTMedTotC \\
    más de 0,75    & \cFTND & \cFTMedD & \cFTMedTotD \\
    \bottomrule
  \end{tabular}
\end{table}

\paragraph{Sensibilidad a la versión de la cobertura.} Se repitieron los análisis con la
cobertura sobre la imagen completa, cuyo denominador no depende de la anotación. Esa versión
sí depende fuertemente de la fracción etiquetada ($r = \cDepFTPearson$), como es inevitable:
una escena con poca superficie etiquetada no puede tener mucha roca etiquetada sobre el total.
Pese a ello, las conclusiones se mantienen. Las dos versiones ordenan las escenas de forma muy
parecida (correlación de rangos de \cSensRho); la diferencia entre máscaras colaborativas y de
experto persiste, con medianas del \cSensTotColab{} y del \cSensTotExp{} en las escenas con
roca ($p \cSensTotP$ en la prueba de Mann-Whitney); y la ausencia de asociación con el conteo
se reproduce ($r = \cSensConteoPearson$, $p = \cSensConteoPearsonP$; $\rho =
\cSensConteoSpearman$, $p = \cSensConteoSpearmanP$).

\subsection{Conteo aproximado de rocas individuales (E2)}

El conteo se aplicó a las \cFlagOk{} escenas de la población de E2 y arrojó un total de
\cNRocas{} rocas, con una mediana de una roca por imagen y un máximo de \cRocasMax. La
Tabla~\ref{tab:bandas-conteo} y la Figura~\ref{fig:bandas-conteo} recogen la distribución.

\begin{table}[H]
  \centering
  \caption[Distribución del conteo de rocas]{Distribución del conteo de rocas sobre las
    \cFlagOk{} escenas de la población de E2.}
  \label{tab:bandas-conteo}
  \begin{tabular}{lrr}
    \toprule
    \textbf{Rocas por imagen} & \textbf{Imágenes} & \textbf{\%} \\
    \midrule
    0 (descartadas por los filtros) & \cBandaCero        & \cBandaCeroPct \\
    1                               & \cBandaUna         & \cBandaUnaPct \\
    2--3                            & \cBandaDosTres     & \cBandaDosTresPct \\
    4--9                            & \cBandaCuatroNueve & \cBandaCuatroNuevePct \\
    10 o más                        & \cBandaDiez        & \cBandaDiezPct \\
    \midrule
    Total de rocas                  & \multicolumn{2}{r}{\cNRocas} \\
    \bottomrule
  \end{tabular}
\end{table}

\begin{figure}[H]
  \centering
  \includegraphics[width=0.82\textwidth]{Figura_15_conteo_por_bandas.png}
  \caption{Distribución del conteo de rocas por imagen.}
  \label{fig:bandas-conteo}
\end{figure}

En \cBandaCero{} escenas —el \cBandaCeroPct\,\% de la población— existen píxeles de roca
grande pero ninguna región supera los filtros de área mínima y forma. La
subsección~\ref{sec:diagnostico} examina esas escenas; el dato se reporta explícitamente
porque delimita el alcance real del indicador.

En promedio, cada escena contiene \cRawMedio{} componentes conectadas antes de la división de
aguas y \cRocasMedio{} rocas después. El paso de división de aguas \textbf{no aumenta} el
conteo promedio, porque el aumento que produce al separar regiones unidas queda compensado por
las componentes que los filtros descartan. Es coherente con que el criterio de prominencia
evitara la fragmentación masiva que producía la detección de máximos por separación mínima
(Sección~\ref{sec:calibracion}), aunque, como se verá, no la eliminó por completo.

\subsubsection{Distribución tamaño--frecuencia}

Clasificando las \cNRocas{} rocas según su tamaño relativo al área de la imagen
(Tabla~\ref{tab:tamanos}, Figura~\ref{fig:tamano-frecuencia}) se obtiene una
\textbf{distribución decreciente}: predominan las rocas pequeñas y la frecuencia disminuye
al aumentar el tamaño.

\begin{table}[H]
  \centering
  \caption{Distribución tamaño--frecuencia de las rocas contadas.}
  \label{tab:tamanos}
  \begin{tabular}{lrr}
    \toprule
    \textbf{Clase de tamaño} & \textbf{Rocas} & \textbf{\%} \\
    \midrule
    Pequeña (\textless\ \SI{0.5}{\percent} del área) & \cTamPeq & \cTamPeqPct \\
    Mediana (\SI{0.5}{\percent} -- \SI{2}{\percent})  & \cTamMed & \cTamMedPct \\
    Grande ($\geq$ \SI{2}{\percent})                  & \cTamGra & \cTamGraPct \\
    \bottomrule
  \end{tabular}
\end{table}

\begin{figure}[H]
  \centering
  \includegraphics[width=0.82\textwidth]{Figura_16_tamano_frecuencia.png}
  \caption{Distribución tamaño--frecuencia de las rocas contadas.}
  \label{fig:tamano-frecuencia}
\end{figure}

Esta forma coincide cualitativamente con las distribuciones tamaño--frecuencia descritas en
los estudios de abundancia de rocas en sitios de aterrizaje
\citep{golombek2003mer, lorenz2023analytic}. La coincidencia debe leerse con cautela: aquí los
tamaños son \textbf{relativos al campo de visión} y no métricos, porque el procedimiento no
dispone de la escala de la escena. La comparación es de forma, no de magnitud.

La solidez media de las regiones aceptadas es de \cSolidezMedia, es decir, formas compactas;
de las \cNConRocas{} escenas con al menos una roca, solo \cNSolidezBaja{} presentan solidez
media inferior a $0{,}7$, señal de contornos cóncavos o de regiones posiblemente subdivididas
en exceso.

\subsection{Relación entre cobertura y conteo (H4)}
\label{sec:dependencia}

La hipótesis H4 sostiene que los dos indicadores aportan información distinta. Un coeficiente
de Pearson próximo a cero no basta para afirmarlo, porque solo excluye la asociación lineal.
La Tabla~\ref{tab:dependencia} reúne las medidas de dependencia descritas en la
Sección~\ref{sec:evaluacion} para este par y para el control del denominador.

\begin{table}[H]
  \centering
  \small
  \caption[Medidas de dependencia]{Medidas de dependencia entre cobertura y conteo, y entre
    fracción etiquetada y cobertura. Intervalos por \textit{bootstrap} y valores $p$ por
    permutación; la información mutua se compara con el percentil 95 de su distribución
    bajo independencia.}
  \label{tab:dependencia}
  \begin{tabular}{llcc}
    \toprule
    \textbf{Medida} & \textbf{Qué detecta} & \textbf{Cobertura--conteo} & \textbf{Fracción--cobertura} \\
     & & ($n = \cDepCCN$) & ($n = \cDepFCN$) \\
    \midrule
    Pearson $r$       & lineal    & \cDepCCPearson{} \cDepCCPearsonIC & \cDepFCPearson{} \cDepFCPearsonIC \\
                      &           & $p = \cDepCCPearsonP$ & $p = \cDepFCPearsonP$ \\
    Spearman $\rho$   & monótona  & \cDepCCSpearman{} \cDepCCSpearmanIC & \cDepFCSpearman{} \cDepFCSpearmanIC \\
                      &           & $p = \cDepCCSpearmanP$ & $p = \cDepFCSpearmanP$ \\
    Kendall $\tau$    & monótona  & \cDepCCKendall{} ($p = \cDepCCKendallP$) & \cDepFCKendall{} ($p = \cDepFCKendallP$) \\
    Correlación de distancia & cualquiera & \cDepCCDcor{} ($p = \cDepCCDcorP$) & \cDepFCDcor{} ($p = \cDepFCDcorP$)$^{a}$ \\
    Información mutua (bits) & cualquiera & \cDepCCIM{} (nula \cDepCCIMNula) & \cDepFCIM{} (nula \cDepFCIMNula) \\
                      &           & $p = \cDepCCIMP$ & $p = \cDepFCIMP$ \\
    \bottomrule
    \multicolumn{4}{l}{\footnotesize $^{a}$ Sobre una submuestra aleatoria de \cDepDcorSub{}
      escenas, por coste de cómputo.}
  \end{tabular}
\end{table}

Los indicadores presentan una \textbf{asociación lineal prácticamente nula} en las escenas
analizadas ($r = \cDepCCPearson$), y tampoco muestran asociación monótona ($\rho =
\cDepCCSpearman$, $p = \cDepCCSpearmanP$). Sin embargo, \textbf{no son estadísticamente
independientes}: la correlación de distancia ($\cDepCCDcor$, $p = \cDepCCDcorP$) y la
información mutua ($p = \cDepCCIMP$) detectan una dependencia débil que los coeficientes de
correlación no capturan. La prueba chi-cuadrado sobre la tabla de contingencia por tramos, más
gruesa, queda en el límite ($\chi^2 = \cChiDos$ con \cChiGl{} grados de libertad,
$p = \cChiP$; V de Cramér \cCramer).

La Tabla~\ref{tab:conteo-por-cobertura} muestra la forma de esa dependencia: una \textbf{U
invertida}. El conteo medio es bajo en las escenas de cobertura muy baja, alcanza su máximo
entre el \SI{10}{\percent} y el \SI{50}{\percent} de cobertura y vuelve a bajar en las escenas
de cobertura total. Los dos extremos tienen explicaciones distintas. En el inferior la relación
es en parte de construcción: los píxeles de roca grande forman parte del numerador de la
cobertura, de modo que una cobertura muy baja implica una región de roca grande pequeña, que
los filtros de área descartan con frecuencia. En el superior, una escena enteramente cubierta
de roca suele ser un afloramiento continuo, que el anotador marca como lecho rocoso y no como
bloques discretos.

\begin{table}[H]
  \centering
  \caption[Conteo según el tramo de cobertura]{Conteo de rocas según el tramo de cobertura,
    sobre la población de E2.}
  \label{tab:conteo-por-cobertura}
  \begin{tabular}{lrccc}
    \toprule
    \textbf{Cobertura} & \textbf{Escenas} & \textbf{Conteo medio} & \textbf{Sin rocas (\%)}
      & \textbf{Cuatro o más (\%)} \\
    \midrule
    hasta \SI{10}{\percent}              & \cUNA & \cUMediaA & \cUCeroA & \cUCuatroA \\
    \SI{10}{\percent} -- \SI{50}{\percent} & \cUNB & \cUMediaB & \cUCeroB & \cUCuatroB \\
    \SI{50}{\percent} -- \SI{80}{\percent} & \cUNC & \cUMediaC & \cUCeroC & \cUCuatroC \\
    \SI{80}{\percent} -- \SI{99.99}{\percent} & \cUND & \cUMediaD & \cUCeroD & \cUCuatroD \\
    \SI{100}{\percent}                   & \cUNE & \cUMediaE & \cUCeroE & \cUCuatroE \\
    \bottomrule
  \end{tabular}
\end{table}

La magnitud de la dependencia se mide con la razón de correlación. Los deciles de cobertura
explican el \cEtaConteo{} de la varianza del conteo, frente a un \cEtaConteoNula{} esperable
por azar ($p = \cEtaConteoP$ por permutación); las bandas de conteo explican el \cEtaCob{} de
la varianza de la cobertura, valor que no se distingue del azar ($p = \cEtaCobP$). Conocer uno
de los indicadores informa, por tanto, muy poco sobre el otro: ninguno es función del otro y
los dos describen aspectos distintos del terreno, cuánta roca hay y cuántos bloques discretos
la componen. Pero la relación existe, no es lineal, y en parte procede de cómo se definen.

\subsection{Composición del terreno y tipología de escenas (E3)}

Además de los dos indicadores previstos, cada imagen se caracterizó por la proporción de
cada clase sobre sus píxeles etiquetados. La composición media es de \cCompBed{} de lecho
rocoso, \cCompSoil{} de suelo, \cCompSand{} de arena y \cCompBig{} de roca grande. La clase
dominante es el lecho rocoso en \cDomBed{} imágenes, el suelo en \cDomSoil, la arena en
\cDomSand{} y la roca grande en apenas \cDomBig.

A partir de esa composición se definió la tipología de la Tabla~\ref{tab:tipologia}, que la
Figura~\ref{fig:tipologia} representa.

\begin{table}[H]
  \centering
  \caption{Tipología de escenas según la composición del terreno.}
  \label{tab:tipologia}
  \begin{tabular}{llrr}
    \toprule
    \textbf{Tipo} & \textbf{Criterio} & \textbf{Imágenes} & \textbf{\%} \\
    \midrule
    Rocoso  & roca $\geq$ \SI{66}{\percent}   & \cTipoRocoso  & \cTipoRocosoPct \\
    Suelo   & suelo $\geq$ \SI{50}{\percent}  & \cTipoSuelo   & \cTipoSueloPct \\
    Arenoso & arena $\geq$ \SI{50}{\percent}  & \cTipoArenoso & \cTipoArenosoPct \\
    Mixto   & ninguna clase predomina         & \cTipoMixto   & \cTipoMixtoPct \\
    Sin etiqueta & sin píxeles clasificados    & \cTipoSinEtiq & \cTipoSinEtiqPct \\
    \bottomrule
  \end{tabular}
\end{table}

\begin{figure}[H]
  \centering
  \includegraphics[width=0.82\textwidth]{Figura_17_tipologia_escenas.png}
  \caption{Tipología de escenas según la composición del terreno.}
  \label{fig:tipologia}
\end{figure}

\subsection{Variación a lo largo de la secuencia de adquisición (E3)}

El orden temporal se obtiene del \textbf{reloj de nave}, el contador de a bordo que el
identificador de cada imagen incluye tras el prefijo de cámara y que se registra en la columna
\texttt{sclk} del conjunto de resultados. Ordenando las imágenes por ese contador y
agrupándolas en tramos de igual número de imágenes, se observa (Figura~\ref{fig:secuencia})
una \textbf{alternancia clara entre segmentos de la secuencia francamente rocosos y segmentos
de suelo o arena}, con concentraciones puntuales de roca grande que alcanzan hasta el
\SI{40}{\percent} de las imágenes de un segmento.

\begin{figure}[H]
  \centering
  \includegraphics[width=\textwidth]{Figura_18_variacion_secuencia.png}
  \caption[Variación a lo largo de la secuencia de adquisición]{Cobertura y presencia de roca
    grande a lo largo de la secuencia de adquisición, en segmentos de igual número de
    imágenes ordenados por el reloj de nave.}
  \label{fig:secuencia}
\end{figure}

Esta lectura debe interpretarse como una \textbf{ordenación relativa en el tiempo} y no como
una variación espacial. El orden cronológico de las imágenes no equivale a distancia recorrida:
el rover puede tomar muchas imágenes desde una misma ubicación, con orientaciones distintas, y
el conjunto no incluye la posición desde la que se adquirió cada una. Por ello el análisis no
permite afirmar en qué lugares del trayecto se concentra la roca, solo en qué momentos de la
secuencia de adquisición.

Cabe señalar que el subconjunto es casi enteramente de la cámara izquierda del par estéreo
—\cNOjoIzq{} imágenes frente a \cNOjoDer{} de la derecha—, por lo que no fue posible la
comparación entre ojos que E3 contemplaba.

\subsection{Extensión aplicada: reglas heurísticas de priorización}
\label{sec:extension}

Como extensión aplicada, y fuera del núcleo de la pregunta de investigación, los indicadores se
tradujeron en seis reglas heurísticas de priorización con umbrales explícitos, y se construyó
una aplicación de escritorio para consultar el conjunto de resultados. Aplicadas al conjunto,
las reglas sitúan \cPrioAlta{} escenas en prioridad alta, \cPrioMedia{} en media y
\cPrioBaja{} en baja; las \cPrioSin{} restantes no activan ninguna regla de prioridad.

Estas reglas \textbf{no constituyen una evaluación de riesgo para la navegación} ni están
calibradas para ella. Sin escala métrica no puede inferirse la altura de un bloque, la
transitabilidad, el daño potencial a las ruedas ni la probabilidad de atrapamiento a partir de
porcentajes de píxeles, y no existe para este subconjunto un registro de incidentes contra el
que validarlas. Su función es ordenar escenas para revisión con un criterio visible y
ajustable. El Anexo~\ref{app:extension} documenta las reglas, su distribución y la aplicación.

\section{Evaluación y validación}

\subsection{Contraste con las máscaras de experto (H3)}
\label{sec:experto}

El conjunto incluye \cNGold{} máscaras etiquetadas por especialistas con acuerdo del
\SI{100}{\percent}, sobre imágenes que \textbf{no} forman parte del conjunto de entrenamiento.
Se ejecutó sobre ellas el mismo procedimiento, sin modificar ningún parámetro
(Tabla~\ref{tab:validacion}, Figura~\ref{fig:validacion}).

\begin{table}[H]
  \centering
  \caption[Indicadores según la fuente de etiqueta]{Indicadores según la fuente de etiqueta.
    Las dos columnas corresponden a imágenes distintas.}
  \label{tab:validacion}
  \begin{tabular}{lcc}
    \toprule
    \textbf{Indicador} & \textbf{Colaborativas} & \textbf{Experto} \\
     & ($n = \cNEscenas$) & ($n = \cNGold$) \\
    \midrule
    Cobertura mediana (escenas con roca)          & \cCobMedVal     & \cGoldCobMedPos \\
    Escenas con cobertura del \SI{100}{\percent}  & \cColabPctCien  & \cGoldPctCien \\
    Lecho rocoso (media sobre lo etiquetado)      & \cColabBed      & \cGoldBed \\
    Suelo y arena (media sobre lo etiquetado)     & \cColabSueloArena & \cGoldSueloArena \\
    Escenas con roca grande                       & \cPctBigAny     & \cGoldUnoBig \\
    Fracción etiquetada (mediana)                 & \cFracValidMed  & \cFracValidGold \\
    \bottomrule
  \end{tabular}
\end{table}

\begin{figure}[H]
  \centering
  \includegraphics[width=0.85\textwidth]{Figura_19_validacion_experto.png}
  \caption[Cobertura según la fuente de etiqueta]{Cobertura según la fuente de etiqueta, en
    las escenas con roca. Las dos distribuciones proceden de muestras de imágenes distintas.}
  \label{fig:validacion}
\end{figure}

\paragraph{Resultado.} Las dos fuentes de anotación presentan distribuciones muy diferentes de
cobertura y de clases. La diferencia de medianas en las escenas con roca es de
\cFuenteDifMed{} puntos porcentuales (IC~95\,\% \cFuenteDifIC; $p \cFuenteMWp$ en la prueba de
Mann-Whitney). En cambio, la fracción etiquetada es prácticamente igual en ambas fuentes
(\cFracValidMed{} frente a \cFracValidGold, $p = \cFracValidP$).

Esta comparación no permite, por sí sola, atribuir la diferencia a la anotación. Las máscaras de
experto cubren \textbf{otras imágenes}, para las que AI4Mars no distribuye máscara
colaborativa, de modo que no es posible comparar las dos fuentes sobre las mismas escenas. Una
diferencia de cobertura puede deberse a que cada grupo anote de forma distinta o a que las dos
muestras contengan terreno distinto.

\paragraph{Instrumento común.} Para separar ambos efectos se aplicó el segmentador entrenado,
que es una función fija de la imagen, a la región anotable de \cFANColab{} escenas
colaborativas no vistas en su entrenamiento y de las \cFANExp{} de experto, según el diseño de
la Sección~\ref{sec:evaluacion}. La Tabla~\ref{tab:instrumento} y la
Figura~\ref{fig:instrumento} presentan el resultado.

\begin{table}[H]
  \centering
  \caption[Descomposición de la diferencia entre fuentes]{Cobertura según la etiqueta y según
    el instrumento común, en las dos muestras de imágenes, y descomposición de la diferencia
    de medias. Intervalos por \textit{bootstrap}.}
  \label{tab:instrumento}
  \small
  \begin{tabular}{lcc}
    \toprule
     & \textbf{Muestra colaborativa} & \textbf{Muestra de experto} \\
     & ($n = \cFANColab$) & ($n = \cFANExp$) \\
    \midrule
    Cobertura según la etiqueta, media (mediana)        & \cFAEtiqMediaColab{} (\cFAEtiqMedColab) & \cFAEtiqMediaExp{} (\cFAEtiqMedExp) \\
    Cobertura según el instrumento, media (mediana)     & \cFARegMediaColab{} (\cFARegMedColab)   & \cFARegMediaExp{} (\cFARegMedExp) \\
    Etiqueta menos instrumento, media                   & \cResColab{} p.p. & \cResExp{} p.p. \\
    \midrule
    \multicolumn{3}{l}{\textit{Descomposición de la diferencia de medias entre muestras}} \\
    Diferencia total según la etiqueta  & \multicolumn{2}{c}{\cDescTotal{} p.p.} \\
    Atribuible a las imágenes           & \multicolumn{2}{c}{\cDescImg{} p.p. \cDescImgIC{} (\cDescImgPct)} \\
    Atribuible a la anotación           & \multicolumn{2}{c}{\cDescAnot{} p.p. \cDescAnotIC{} (\cDescAnotPct)} \\
    \bottomrule
  \end{tabular}
\end{table}

\begin{figure}[H]
  \centering
  \includegraphics[width=\textwidth]{Figura_30_instrumento_comun.png}
  \caption[Instrumento común sobre las dos muestras]{Izquierda: distribución acumulada de la
    cobertura que el segmentador asigna a la región anotable de cada escena, sin intervención
    de ninguna etiqueta. Derecha: descomposición de la diferencia de cobertura media entre las
    dos muestras, con intervalos del \SI{95}{\percent}.}
  \label{fig:instrumento}
\end{figure}

El instrumento común reproduce la mayor parte de la diferencia sin pasar por ninguna anotación.
Medidas por el mismo segmentador, las imágenes de la muestra de experto tienen una cobertura
mediana del \cFARegMedExp{} frente al \cFARegMedColab{} de la muestra colaborativa: \textbf{el
conjunto de máscaras de experto está compuesto por escenas mucho menos rocosas}. De la
diferencia total de medias, \cDescImg{} puntos —el \cDescImgPct— se asocian con la composición
de las muestras de imágenes y \cDescAnot{} puntos —el \cDescAnotPct— con la fuente de las
etiquetas. Esta segunda componente es pequeña y su intervalo apenas excluye el cero. Ambas
anotaciones asignan algo menos de roca que el instrumento, la colaborativa \cResColab{} puntos y
la de experto \cResExp: la diferencia entre ambos residuos es la componente de anotación.

\paragraph{Sensibilidad.} La conclusión no depende de la muestra ni del instrumento concretos.
Tomando la muestra colaborativa solo de los bloques temporales de prueba ($n = \cFPNColab$), a más
de un sol de cualquier imagen de entrenamiento, la componente de imágenes es el \cFPDescImgPct{}
y la de anotación \cFPDescAnot{} puntos \cFPDescAnotIC. Con el segmentador de una versión
anterior, entrenado con un reparto al azar por imagen, la descomposición atribuía a las imágenes
el \cFVDescImgPct{} y a la anotación \cFVDescAnot{} puntos \cFVDescAnotIC. En todos los casos la
mayor parte de la diferencia procede de las imágenes y la componente de anotación es pequeña y
positiva; su magnitud exacta, entre dos y cuatro puntos, depende del instrumento.

Las imágenes de experto no proceden de un periodo distinto de la misión: la mediana de la
distancia, en reloj de nave, a la imagen de entrenamiento más próxima es de \cFugaExpSoles{}
soles. La diferencia de composición refleja, por tanto, cómo se seleccionaron las escenas del
conjunto de prueba, algo que la documentación del conjunto no especifica y este trabajo no
puede determinar.

La descomposición tiene un supuesto que conviene explicitar: que el instrumento responde igual
a las dos muestras. Como el segmentador se entrenó con etiquetas colaborativas, puede haber
heredado parte del criterio de esa anotación; si así fuera, aplicado a las imágenes de experto
tendería a asignarles más roca de la que el experto marca, y la componente atribuida a la
anotación quedaría subestimada. Es lo que se observa: frente al experto, el modelo sobreestima
la cobertura en \cBAExpMedia{} puntos de media (Sección~\ref{sec:modelo-experto}). La cifra de
\cDescAnot{} puntos debe leerse, por tanto, como una estimación aproximada y probablemente
conservadora.

\paragraph{Interpretación.} Dado que el mismo cálculo se aplicó a ambas fuentes, la diferencia
observada no puede atribuirse a la implementación de la fórmula; según el instrumento común, se
asocia principalmente con la composición de las muestras de imágenes y, en menor medida, con la
fuente de las etiquetas. La componente de anotación es compatible con la hipótesis de un
\textbf{sesgo de saliencia} en la anotación colaborativa —que las regiones visualmente
prominentes, como la roca, reciban una anotación más consistente o completa—, pero el diseño
del estudio no permite identificar causalmente ese mecanismo: no se midieron la atención, la
fatiga ni el tiempo de anotación. Un mecanismo alternativo que se había planteado, que la
anotación colaborativa dejara sin etiquetar más suelo y arena y eso redujera el denominador,
\textbf{no encuentra apoyo}: la fracción etiquetada es igual en las dos fuentes.

\subsection{Validación con conteo humano (H2)}
\label{sec:validacion-manual}

La metodología comprometía un contraste entre el conteo automático y un conteo manual, con el
diseño descrito en la Sección~\ref{sec:evaluacion}. Se realizó en dos rondas: una piloto de
\num{24} escenas y una definitiva de \cVN{} escenas, que es la que se reporta. Los parámetros
estaban congelados antes de la primera respuesta y ninguna escena de calibración forma parte de
la muestra definitiva (Sección~\ref{sec:calibracion}). La validación la realizó un único
observador, de modo que lo que sigue es \textbf{una evaluación exploratoria} y no una medición
de exactitud.

\paragraph{Correcciones derivadas de la ronda piloto.} La ronda piloto empleaba cuatro bandas,
la última abierta en «diez o más». Ese diseño resultó defectuoso: un procedimiento que contara
ochenta rocas donde el observador distinguía una decena puntuaba como acierto exacto, de modo
que la escala \textbf{favorecía sistemáticamente a los métodos que sobreestiman}. La ronda
definitiva subdivide esa banda. Corrige además otros tres defectos: excluye las escenas cuya
región anotada no alcanza los \num{2000} píxeles —en la ronda piloto siete de veinticuatro eran
imperceptibles—, sustituye el relleno opaco por un tinte tenue que no oculta la textura, y
estratifica la muestra por área de la región en lugar de por la banda del propio
procedimiento.

\subsubsection{Resultado}

Sobre las \cVN{} escenas, el acuerdo exacto de banda es del \cVAcuerdo, con coeficiente kappa
de Cohen de \cVKappa{} (IC~95\,\% \cVKappaIC) y kappa ponderado de \cVKappaPond{} (IC~95\,\%
\cVKappaPondIC). Ambos corresponden a un acuerdo \textbf{escaso} en la escala de
\citet{landis1977}. El acuerdo aparente es engañoso: \cVCoincUnaTres{} de las \cVCoinc{}
coincidencias se concentran en una sola celda, la de una a tres rocas.

El desacuerdo es \textbf{marcadamente asimétrico}. El procedimiento sitúa la escena en una
banda inferior a la del observador en \cVAbajo{} casos y superior en solo \cVArriba{}
($p \cVSignoP$ en la prueba de signos). La Figura~\ref{fig:validacion-manual} muestra la
matriz de acuerdo.

\begin{figure}[H]
  \centering
  \includegraphics[width=0.78\textwidth]{Figura_26_validacion_manual_v2.png}
  \caption[Acuerdo entre el conteo manual y el automático]{Matriz de acuerdo entre el
    conteo manual por bandas y el automático, sobre \cVN{} escenas. Las tres columnas de
    la derecha están vacías: con los parámetros fijados, el procedimiento no produce ningún
    conteo superior a nueve.}
  \label{fig:validacion-manual}
\end{figure}

El rasgo más visible de esa matriz son las \textbf{tres columnas vacías}. En ninguna de las
\cVN{} escenas el procedimiento devuelve más de \cVMaxAuto{} rocas, mientras que el observador
identificó diez o más en \cVHumDiez{} de ellas y veinticinco o más en \cVHumVeinti. Sumando las
escenas, el procedimiento cuenta \cVRocasAuto{} rocas donde el observador distingue al menos
\cVRocasHumMin{} —el límite inferior de sus bandas—, es decir, \cVRazon{} veces más.

\paragraph{Sensibilidad a los parámetros.} Para descartar que el resultado dependa de la
calibración se repitió la comparación con las \cSenN{} combinaciones de prominencia, suavizado
y área mínima descritas en la Sección~\ref{sec:calibracion}. El kappa ponderado oscila entre
\cSenMin{} y \cSenMax, con mediana de \cSenMed; en todas las combinaciones el procedimiento
queda por debajo del observador en entre \cSenAbajoMin{} y \cSenAbajoMax{} escenas, y por
encima en no más de \cSenArribaMax. Solo \cSenDiez{} combinaciones, las más permisivas,
alcanzan alguna vez la banda de diez a veinticuatro rocas. El subconteo no es, por tanto, un
artefacto del ajuste de los umbrales.

\subsubsection{Por qué el subconteo: anotación no exhaustiva y regiones sin información geométrica}
\label{sec:anotacion-parcial}

El examen de las escenas discrepantes identifica dos causas, ninguna de las cuales se corrige
ajustando el procedimiento.

La primera es que la anotación de roca grande \textbf{no marca todas las rocas de la escena,
sino algunas}. En \cVNoExhN{} de las \cVN{} escenas, la clase roca grande representa menos del
\SI{20}{\percent} de la roca etiquetada; el resto se anotó como lecho rocoso. La mediana es del
\cVNoExhMed. En las \cVUnaTresN{} escenas que el observador situó entre una y tres rocas, la
región resaltada ocupa una mediana del \cVUnaTresZona{} de la imagen y del \cVUnaTresFr{} de la
roca etiquetada: regiones diminutas que contienen, en efecto, una o dos rocas, mientras la
escena contiene muchas más fuera de ellas.

La segunda es que algunas regiones \textbf{no contienen la información geométrica necesaria
para separar instancias}. La división de aguas sobre la transformada de distancia solo puede
cortar donde la región presenta un estrangulamiento. Cuando el anotador traza un único polígono
que envuelve un campo de bloques contiguos, la región resultante es convexa o casi convexa, y
la máscara no contiene ningún indicio de dónde termina un bloque y empieza el siguiente. La
Figura~\ref{fig:mecanismo} contrasta un caso en que el procedimiento coincide con el observador
y otro en que, sobre una sola región que el observador cuenta entre veinticinco y cuarenta y
nueve rocas, el procedimiento solo encuentra tres.

\begin{figure}[H]
  \centering
  \includegraphics[width=\textwidth]{Figura_28_mecanismo_subconteo.png}
  \caption[Mecanismo del subconteo]{Arriba, una escena en la que la anotación delimita cada
    bloque por separado y el conteo coincide con el observador. Abajo, un único polígono
    envuelve un campo de bloques: la transformada de distancia no presenta estrangulamientos
    y la división de aguas no tiene dónde cortar.}
  \label{fig:mecanismo}
\end{figure}

\subsubsection{Alcance del indicador}

El procedimiento produce un conteo reproducible de las instancias geométricamente separables
dentro de la clase roca grande, pero ese conteo no puede interpretarse como estimación válida
del número total de rocas visibles en la escena. La evaluación exploratoria con un observador
sugiere un subconteo sistemático, atribuible tanto a la cobertura incompleta de la clase como a
regiones que no contienen la información geométrica necesaria para separar instancias.

Lo que el resultado sostiene es un uso más débil: ordenar de forma aproximada escenas según
contengan o no bloques discretos delimitados, sabiendo que el valor queda por debajo de lo que
un observador contaría. La conclusión debe moderarse además por el diseño de la validación: con
un único observador no es posible separar el desacuerdo entre algoritmo y persona de la
variabilidad natural entre personas, y la muestra, estratificada por área de la región, no es
representativa del subconjunto completo.

\subsection{Conteo con información de la imagen: exploración}
\label{sec:hibrido}

Si la máscara no contiene la información que separa los bloques, esa información solo puede
proceder de la imagen. Se exploró esa vía manteniendo el procedimiento clásico: la región
sigue siendo la que la anotación marcó como roca grande, pero el relieve sobre el que corta la
división de aguas se construye a partir de la imagen. Se compararon cuatro relieves: el
gradiente de intensidad a escala fina; el gradiente a escala gruesa, en el que la textura
interna se atenúa y persiste el borde; un realce de líneas oscuras finas mediante
\textit{top-hat} negro, que corresponde a las sombras entre rocas; y una combinación a partes
iguales de la geometría de la máscara y el gradiente. El protocolo es idéntico para los cuatro
—mismas semillas por prominencia y mismos filtros finales que E2—, y su único parámetro se
eligió por validación cruzada dejando una escena fuera, para que ninguna escena influyera en su
propio resultado. La diferencia con E2 se contrastó por \textit{bootstrap}, primero al
\SI{95}{\percent} y después con corrección de Bonferroni por el número de métodos probados
(Tabla~\ref{tab:metodos-imagen}).

\begin{table}[H]
  \centering
  \small
  \caption[Conteo con relieves derivados de la imagen]{Acuerdo con el observador de los
    conteos que usan la imagen dentro de la región anotada, frente al kappa ponderado de E2
    (\cVKappaPond).}
  \label{tab:metodos-imagen}
  \setlength{\tabcolsep}{4pt}
  \begin{tabular}{lccccc}
    \toprule
     & & & \multicolumn{3}{c}{\textbf{Diferencia de $\kappa$ ponderado con E2}} \\
    \cmidrule(lr){4-6}
    \textbf{Relieve} & \textbf{$\kappa$ pond.} & \textbf{Debajo / encima}
      & \textbf{Media} & \textbf{IC 95\,\%} & \textbf{IC Bonferroni} \\
    \midrule
    Gradiente fino     & \cMIGradKw    & \cMIGradAbajo{} / \cMIGradArriba       & \cMIGradDif    & \cMIGradIC    & \cMIGradICB \\
    Gradiente grueso   & \cMIPersistKw & \cMIPersistAbajo{} / \cMIPersistArriba & \cMIPersistDif & \cMIPersistIC & \cMIPersistICB \\
    Sombras            & \cMISombraKw  & \cMISombraAbajo{} / \cMISombraArriba   & \cMISombraDif  & \cMISombraIC  & \cMISombraICB \\
    Combinado          & \cMIComboKw   & \cMIComboAbajo{} / \cMIComboArriba     & \cMIComboDif   & \cMIComboIC   & \cMIComboICB \\
    \bottomrule
  \end{tabular}
\end{table}

Ningún relieve mejora a E2 de forma demostrable. El de sombras obtiene el mayor acuerdo y su
intervalo al \SI{95}{\percent} excluye el cero, pero deja de hacerlo al corregir por el número
de métodos probados. La sensibilidad lo confirma: sus \cMISombraVarN{} variantes de suavizado y
radio dan kappas ponderados entre \cMISombraVarMin{} y \cMISombraVarMax{}, y ninguna supera la
corrección de Bonferroni. Los cuatro relieves invierten además la dirección del error: todos
quedan por encima del observador en más escenas que por debajo —el de escala gruesa, en
\cMIPersistArriba{} de las \cVN—, lo que indica que cortan también en la textura interna de
las rocas y no solo en sus bordes. Corregir el subconteo con la imagen exige, por tanto,
controlar a la vez la sobresegmentación que la propia imagen introduce.

La Figura~\ref{fig:hibrido} ilustra el principio sobre la escena de la
Figura~\ref{fig:diag-bedrock}, etiquetada casi por completo como lecho rocoso y con conteo
nulo: el gradiente de la imagen, restringido a la región de roca de la máscara
(\cHibridoRegion{} de la escena), delimita \cHibridoN{} bloques cuyas fronteras siguen las
fracturas visibles. La cifra depende del umbral de prominencia de las semillas, es con toda
probabilidad un sobreconteo y no está validada.

\begin{figure}[H]
  \centering
  \includegraphics[width=\textwidth]{diag_prueba_hibrido.png}
  \caption[Prueba de concepto del enfoque híbrido]{Prueba de concepto: la máscara indica
    dónde hay roca y el gradiente de la imagen dónde están las fronteras entre bloques. El
    resultado ilustra el principio y no constituye un conteo validado.}
  \label{fig:hibrido}
\end{figure}

El enfoque híbrido se mantiene, por tanto, como \textbf{prueba de concepto}: muestra que la
imagen contiene información sobre las fronteras entre bloques que la máscara no recoge, pero
la evidencia reunida no permite concluir que la aproveche mejor que E2, y con \cVN{} escenas
y un solo observador la comparación carece de potencia para detectar mejoras moderadas. Su
evaluación con más observadores y una muestra independiente queda como trabajo futuro.

\subsection{Comparación con métodos de aprendizaje automático}

\paragraph{Modelo general sin entrenamiento específico.} El modelo fundacional de segmentación
\citep{kirillov2023sam}, en su variante rápida \citep{zhao2023fastsam}, se aplicó a \cSamN{}
imágenes con roca grande, restringiendo sus regiones a la zona etiquetada como roca. El acuerdo
con el conteo clásico, medido por bandas, es del \cSamAcuerdo, con correlación de rangos de
\cSamRho{} y un error absoluto medio de \cSamMAE{} rocas (Figura~\ref{fig:matriz-acuerdo}). El
modelo tiende a subdividir una misma roca según su textura interna y, en otras escenas, a no
detectar bloques de bajo contraste.

\begin{figure}[H]
  \centering
  \includegraphics[width=0.72\textwidth]{Figura_20_matriz_acuerdo.png}
  \caption{Matriz de acuerdo entre el conteo clásico y el modelo general.}
  \label{fig:matriz-acuerdo}
\end{figure}

\paragraph{Segmentador entrenado con las propias máscaras.} El segmentador DeepLabV3, con la
configuración y el reparto por bloques temporales de la Tabla~\ref{tab:deeplab-config}, alcanzó
su mejor mIoU de validación en la época \cSegEpocaElegida{} (\cSegValMIoU, frente a
\cSegValUltima{} en la última), y ese es el modelo que se evaluó, una sola vez, sobre la prueba.
La Tabla~\ref{tab:deeplab} recoge su desempeño en tres conjuntos: la muestra equilibrada de
prueba, todas las escenas elegibles de los bloques de prueba con su distribución natural, y las
escenas de esos bloques que el umbral de fracción etiquetada había excluido.

\begin{table}[H]
  \centering
  \small
  \caption[Desempeño del segmentador entrenado]{Desempeño del segmentador sobre los bloques
    temporales de prueba, ninguno a menos de un sol de una imagen de entrenamiento. Error en
    puntos porcentuales; con signo, modelo menos etiqueta.}
  \label{tab:deeplab}
  \begin{tabular}{lccc}
    \toprule
     & \textbf{Muestra equilibrada} & \textbf{Distribución natural} & \textbf{Excluidas por fracción} \\
     & ($n = \cSegTest$) & ($n = \cSegNatN$) & ($n = \cSegExcN$) \\
    \midrule
    IoU de la clase no-roca       & \cSegIoUNo    & \cSegNatIoUNo    & \cSegExcIoUNo \\
    IoU de la clase roca          & \cSegIoURoca  & \cSegNatIoURoca  & \cSegExcIoURoca \\
    IoU medio                     & \textbf{\cSegMIoU} & \textbf{\cSegNatMIoU} & \cSegExcMIoU \\
    Correlación de la cobertura   & \cSegRCob     & \cSegNatR        & \cSegExcR \\
    Error absoluto medio          & \cSegMAE      & \cSegNatMAE      & \cSegExcMAE \\
    Error medio con signo         & \cSegErrMedio & \cSegNatErr      & \cSegExcErr \\
    \bottomrule
  \end{tabular}
\end{table}

El desempeño sobre la distribución natural es muy próximo al de la muestra equilibrada, de modo
que el equilibrio artificial de la prueba no infla las cifras. Las escenas con menos del
\SI{20}{\percent} de superficie etiquetada, en cambio, se segmentan claramente peor —error
absoluto medio de \cSegExcMAE{} puntos—: la exclusión no solo protege la función de pérdida,
también delimita el dominio en que el modelo es fiable.

\begin{figure}[H]
  \centering
  \includegraphics[width=0.62\textwidth]{Figura_21_cobertura_modelo_vs_humano.png}
  \caption[Cobertura del modelo frente a la humana]{Cobertura estimada por el modelo
    entrenado frente a la derivada de la anotación colaborativa, sobre las \cSegTest{}
    imágenes de la muestra equilibrada de prueba.}
  \label{fig:modelo-vs-humano}
\end{figure}

\paragraph{Fuga de información.} Con el reparto por bloques y el margen de un sol, la imagen de
prueba más próxima a una de entrenamiento está a \cFugaMinSoles{} soles, y la distancia mediana
es de \cFugaMedGapSoles{} soles. Una versión anterior del modelo, entrenada con un reparto al azar
por imagen, tenía \cAleMinuto{} imágenes de prueba (\cAleMinutoPct) a menos de un minuto de una
de entrenamiento, y obtenía un IoU medio de \cAleMIoU{} y una correlación de la cobertura de
\cAleR. Que el reparto por bloques entregue prácticamente las mismas cifras indica que aquella
proximidad no inflaba el desempeño. La Tabla~\ref{tab:fuga} muestra además que, dentro de la
prueba, el desempeño no depende de la distancia al entrenamiento.

\begin{table}[H]
  \centering
  \caption[Desempeño según la distancia temporal al entrenamiento]{Desempeño del segmentador
    sobre la muestra equilibrada de prueba según la distancia, en reloj de nave, de cada imagen
    a la imagen de entrenamiento más próxima.}
  \label{tab:fuga}
  \begin{tabular}{lrcc}
    \toprule
    \textbf{Distancia al entrenamiento} & \textbf{Imágenes} & \textbf{IoU medio}
      & \textbf{Correlación de la cobertura} \\
    \midrule
    De 1 a 10 soles     & \cFugaNA & \cFugaIoUA & \cFugaRA \\
    De 10 a 50 soles    & \cFugaNB & \cFugaIoUB & \cFugaRB \\
    Más de 50 soles     & \cFugaNC & \cFugaIoUC & \cFugaRC \\
    \bottomrule
  \end{tabular}
\end{table}

\subsubsection{Contraste del modelo con etiquetas de experto}
\label{sec:modelo-experto}

Se evaluó el mismo modelo, sin reentrenarlo, contra las \cNGold{} máscaras de experto,
restringiendo la comparación a los píxeles que el experto etiquetó. El IoU medio desciende a
\cExpMIoU{} (\cExpIoUNo{} en no-roca y \cExpIoURoca{} en roca) y la correlación de la cobertura
a \cExpR.

Una correlación alta no indica, por sí sola, ausencia de sesgo: mide si el orden de las escenas
se preserva, no si los valores coinciden. Para valorar el acuerdo se examinó el error con signo,
definido como la cobertura del modelo menos la de la etiqueta, tanto frente al experto como
frente a la muestra colaborativa aleatoria del apartado anterior
(Tabla~\ref{tab:modelo-experto}, Figuras~\ref{fig:modelo-experto}
y~\ref{fig:bland-altman}).

\begin{table}[H]
  \centering
  \caption[Error de la cobertura del modelo según la referencia]{Error de la cobertura del
    modelo según la referencia, en puntos porcentuales. Intervalos por \textit{bootstrap};
    límites de acuerdo de Bland--Altman como media $\pm\,1{,}96$ desviaciones típicas.}
  \label{tab:modelo-experto}
  \begin{tabular}{lcc}
    \toprule
    \textbf{Medida} & \textbf{Colaborativa} & \textbf{Experto} \\
     & ($n = \cBAColN$) & ($n = \cBAExpN$) \\
    \midrule
    Error medio, IC 95\,\%           & \cBAColMedia{} \cBAColMediaIC & \cBAExpMedia{} \cBAExpMediaIC \\
    Error mediano                    & \cBAColMediana & \cBAExpMediana \\
    Error absoluto medio, IC 95\,\%  & \cBAColMAE{} \cBAColMAEIC & \cBAExpMAE{} \cBAExpMAEIC \\
    Límites de acuerdo               & \cBAColLoA & \cBAExpLoA \\
    Escenas con error mayor de 10 p.p. & \cBAColDiez & \cBAExpDiez \\
    Escenas con error mayor de 50 p.p. & \cBAColCincuenta & \cBAExpCincuenta \\
    \bottomrule
  \end{tabular}
\end{table}

\begin{figure}[H]
  \centering
  \includegraphics[width=\textwidth]{Figura_27_modelo_vs_experto.png}
  \caption[Cobertura del modelo frente a la del experto]{Cobertura estimada por el modelo
    frente a la derivada de la anotación de experto, y distribución del error con signo,
    sobre las \cNGold{} escenas con máscara de especialista.}
  \label{fig:modelo-experto}
\end{figure}

\begin{figure}[H]
  \centering
  \includegraphics[width=\textwidth]{Figura_29_bland_altman.png}
  \caption[Diagrama de Bland--Altman]{Diferencia entre la cobertura del modelo y la de la
    etiqueta frente a su media, con la referencia colaborativa y con la de experto. Línea
    continua: error medio; discontinuas: límites de acuerdo.}
  \label{fig:bland-altman}
\end{figure}

Frente a la anotación colaborativa, el error medio es de \cBAColMedia{} puntos con un intervalo
que incluye el cero \cBAColMediaIC: en agregado, el modelo reproduce la referencia de la que
aprendió. Frente al experto, en cambio, \textbf{sobreestima}: el error medio es de
\cBAExpMedia{} puntos, con un intervalo que excluye el cero \cBAExpMediaIC{}, y el modelo queda
por encima del experto en más escenas que por debajo. Ninguna de las dos cifras autoriza a
llamar al modelo insesgado, y la mediana nula no lo contradice: refleja que en el \cExpExactas{}
de las escenas modelo y etiqueta coinciden exactamente, casi siempre porque ninguno de los dos
encuentra roca. El error absoluto medio frente al experto es de \cBAExpMAE{} puntos, los límites
de acuerdo abarcan unos sesenta y cinco puntos porcentuales y el \cBAExpDiez{} de las escenas
supera los diez puntos de error.

El error tampoco es homogéneo (Tabla~\ref{tab:error-tramos}). Frente al experto, la
sobreestimación se concentra en las coberturas intermedias —entre el \SI{25}{\percent} y el
\SI{75}{\percent} el error medio supera los diez puntos— y se invierte en las escenas de
cobertura total; por tipo de escena, el error es pequeño en las de suelo y mayor en las mixtas.
Un error medio moderado resulta, en parte, de errores de signo opuesto en distintos tramos que
se compensan.

\begin{table}[H]
  \centering
  \small
  \caption[Error del modelo por tramo de cobertura y tipo de escena]{Error de la cobertura
    del modelo por tramo de cobertura de la referencia y por tipo de escena: número de
    escenas, error medio con signo y error absoluto medio, en puntos porcentuales.}
  \label{tab:error-tramos}
  % ARCHIVO GENERADO por scripts/generar_cifras.py — no editar a mano.
\begin{tabular}{lrrrrrr}
\toprule
 & \multicolumn{3}{c}{\textbf{Referencia colaborativa}} & \multicolumn{3}{c}{\textbf{Referencia de experto}} \\
\cmidrule(lr){2-4}\cmidrule(lr){5-7}
 & $n$ & Medio & Absoluto & $n$ & Medio & Absoluto \\
\midrule
\multicolumn{7}{l}{\textit{Por tramo de cobertura de la referencia (\%)}} \\
0 & 171 & \ensuremath{+}\num{1.3} & \num{1.3} & 108 & \ensuremath{+}\num{0.6} & \num{0.6} \\
0--25 & 68 & \ensuremath{+}\num{2.2} & \num{7.2} & 82 & \num{-0.3} & \num{8.3} \\
25--50 & 28 & \ensuremath{+}\num{6.6} & \num{11.6} & 28 & \ensuremath{+}\num{13.0} & \num{23.6} \\
50--75 & 49 & \ensuremath{+}\num{3.4} & \num{13.6} & 28 & \ensuremath{+}\num{16.3} & \num{19.8} \\
75--99,99 & 97 & \ensuremath{+}\num{1.7} & \num{4.8} & 55 & \ensuremath{+}\num{5.1} & \num{11.7} \\
100 & 180 & \num{-1.9} & \num{1.9} & 21 & \num{-6.9} & \num{6.9} \\
\midrule
\multicolumn{7}{l}{\textit{Por tipo de escena}} \\
Suelo & 200 & \ensuremath{+}\num{1.3} & \num{2.9} & 142 & \num{-1.5} & \num{3.7} \\
Arenoso & 64 & \ensuremath{+}\num{5.1} & \num{7.0} & 69 & \ensuremath{+}\num{6.1} & \num{9.5} \\
Rocoso & 298 & \num{-0.4} & \num{3.7} & 89 & \ensuremath{+}\num{4.9} & \num{12.2} \\
Mixto & 31 & \ensuremath{+}\num{2.7} & \num{13.1} & 22 & \ensuremath{+}\num{16.5} & \num{22.3} \\
\bottomrule
\end{tabular}

\end{table}

\paragraph{Lectura.} El segmentador estima la cobertura con un desempeño alto frente a la
anotación de la que aprendió y algo menor frente al experto, respecto del cual sobreestima unos
tres puntos de media, con errores individuales grandes y dependientes del tramo. Esa
sobreestimación es la esperable si el modelo heredó el criterio de la anotación colaborativa. Sirve, por tanto, para describir conjuntos de escenas y no para valorar una escena
concreta. Estima cobertura, no conteo: distingue roca de no-roca y no separa bloques, de modo
que no sustituye a E2. Con esas reservas, que la imagen sola permita aproximar la cobertura
sugiere que el indicador podría calcularse sin anotación humana, posibilidad que el
Capítulo~\ref{cap:conclusiones} recoge como línea futura.

\section{Implicaciones y limitaciones}

\subsection{Contraste de las hipótesis}

La Tabla~\ref{tab:contraste-hipotesis} resume el contraste de las hipótesis de la
Sección~\ref{sec:hipotesis} con la evidencia de este capítulo.

\begin{table}[H]
  \centering
  \small
  \caption{Contraste de las hipótesis de trabajo.}
  \label{tab:contraste-hipotesis}
  \begin{tabular}{p{1.0cm}p{8.6cm}p{3.6cm}}
    \toprule
    \textbf{Hip.} & \textbf{Evidencia} & \textbf{Dictamen} \\
    \midrule
    H1 & El cálculo es determinista y reproducible. Sin asociación lineal con la fracción
      etiquetada ($r = \cDepFCPearson$); dependencia no lineal débil pero detectable
      (correlación de distancia \cDepFCDcor); la cobertura por tramos de fracción no sigue el
      patrón que produciría el artefacto del denominador.
      & Se sostiene, con dependencia residual débil \\
    H2 & Kappa ponderado \cVKappaPond{} \cVKappaPondIC; \cVAbajo{} escenas por debajo del
      observador y \cVArriba{} por encima; subconteo en las \cSenN{} combinaciones de
      parámetros.
      & Se rechaza, en una evaluación exploratoria \\
    H3 & Diferencia de medianas de \cFuenteDifMed{} puntos entre fuentes. El instrumento común
      atribuye el \cDescImgPct{} de la diferencia de medias a las imágenes y \cDescAnot{}
      puntos \cDescAnotIC{} a la anotación.
      & Se rechaza: hay efecto de la fuente, menor que la diferencia bruta \\
    H4 & Sin asociación lineal ni monótona ($r = \cDepCCPearson$, $\rho = \cDepCCSpearman$);
      dependencia no lineal débil (correlación de distancia \cDepCCDcor, $p = \cDepCCDcorP$)
      en forma de U invertida; los deciles de cobertura explican el \cEtaConteo{} de la
      varianza del conteo.
      & Se sostiene: información distinta, no independencia \\
    \bottomrule
  \end{tabular}
\end{table}

\subsection{Interpretación de los indicadores}

El trabajo confirma que las máscaras de AI4Mars pueden leerse como mapas cuantitativos y no
solo como insumo de entrenamiento, con límites distintos para cada indicador.

La cobertura es el indicador sólido. Se calcula para el \cPctCobCalc{} de las escenas, es
reproducible y no está gobernada por la fracción de escena etiquetada. Su valor, sin embargo,
depende de quién anotó: aplicada a máscaras de experto produciría coberturas algo menores, y
la comparación directa con esas máscaras está dominada por el hecho de que cubren escenas
mucho menos rocosas. Las coberturas reportadas describen, por tanto, el terreno tal como lo
registra la anotación colaborativa, y son adecuadas para comparar escenas del mismo conjunto.

El conteo es el indicador limitado. Cuenta de forma reproducible las instancias que la
geometría de la máscara permite separar, pero no el número de rocas que un observador
distingue. Cobertura y conteo aportan información distinta —conocer uno informa muy poco sobre
el otro—, de modo que reportarlos por separado está justificado, aunque no sean
estadísticamente independientes.

\subsection{Diagnóstico de los modos de error del conteo}
\label{sec:diagnostico}

Al inspeccionar visualmente los resultados se observaron dos comportamientos opuestos:
escenas en las que el conteo señalaba rocas donde un observador no las percibía, y escenas con
bloques evidentes en las que el conteo era nulo. Puesto que una explicación natural sería la
insuficiencia de las imágenes, se hizo un diagnóstico sistemático para atribuir cada modo de
error a su causa.

La primera constatación es de orden metodológico: \textbf{el procedimiento no lee las
imágenes}, sino las máscaras, de modo que su calidad no interviene directamente en el conteo.
Las imágenes son además de \num{1024}$\times$\num{1024} píxeles, resolución completa de la
cámara, y las máscaras tienen exactamente el mismo tamaño, por lo que tampoco hay pérdida por
remuestreo.

La Tabla~\ref{tab:diagnostico} reparte las \cNEscenas{} escenas según el resultado del
conteo.

\begin{table}[H]
  \centering
  \caption{Reparto de las escenas según el resultado del conteo.}
  \label{tab:diagnostico}
  \small
  \begin{tabular}{lrrl}
    \toprule
    \textbf{Situación} & \textbf{Escenas} & \textbf{\%} & \textbf{Naturaleza} \\
    \midrule
    La máscara no tiene ningún píxel de roca grande & \cNSinBig     & \cPctSinBig    & límite de la etiqueta \\
    Roca grande en escena casi sin etiquetar (excluida) & \cNNullBig & \cPctNullBig & población de E2 \\
    Roca grande, pero los filtros la descartan       & \cNFiltradas  & \cPctFiltradas & límite de los umbrales \\
    Al menos una roca contada                        & \cNConRocas   & \cPctConRocas  & --- \\
    \bottomrule
  \end{tabular}
\end{table}

El primer grupo es con mucho el mayor, y su componente principal son las \cFlagNoBig{} escenas
con bandera \texttt{no\_bigrock}: \textbf{hay roca visible, pero el anotador la etiquetó como
lecho rocoso}. La mediana de lecho rocoso en ese grupo es del \cNoBigBedMed, y
\cNoBigBedOchenta{} escenas tienen más del \SI{80}{\percent} de lecho rocoso y ninguna roca
grande. El desequilibrio es estructural: sobre las escenas con roca, el lecho rocoso aporta el
\cParteBed{} de la roca etiquetada y la roca grande solo el \cParteBig, y E2 cuenta únicamente
esta última.

La Figura~\ref{fig:diag-bedrock} muestra el caso con claridad: una escena con bloques
individuales distinguibles, etiquetada en su totalidad como lecho rocoso, sobre la que el
conteo devuelve cero.

\begin{figure}[H]
  \centering
  \includegraphics[width=\textwidth]{diag_bedrock_no_contado.png}
  \caption[Roca visible etiquetada como lecho rocoso]{Escena con bloques individuales
    evidentes, etiquetada íntegramente como lecho rocoso. El conteo devuelve cero porque
    la clase roca grande no está presente en la máscara.}
  \label{fig:diag-bedrock}
\end{figure}

El grupo de escenas cuyos píxeles de roca grande descartan los filtros —\cNFiltradas, el
\cPctFiltradas{} del conjunto— se revisó y los filtros \textbf{cumplen su función}: se trata de
motas de pocos píxeles, con un ejemplo de una única región de doce píxeles frente al umbral de
\num{524}, o de bandas muy alargadas trazadas sobre el horizonte, descartadas por relación de
aspecto (Anexo~\ref{app:material-complementario}). No se encontraron motivos para relajar los
umbrales.

Por el lado del conteo excesivo, la Tabla~\ref{tab:falsos-positivos} recoge las tres fuentes
identificadas.

\begin{table}[H]
  \centering
  \caption{Fuentes de conteo excesivo.}
  \label{tab:falsos-positivos}
  \begin{tabular}{lrl}
    \toprule
    \textbf{Fuente} & \textbf{Magnitud} & \textbf{Naturaleza} \\
    \midrule
    Subdivisión por división de aguas & \cNSubdiv{} de \cNConRocas{} escenas & algoritmo \\
    Artefacto de anotación            & \cNArtefacto{} escenas, todas excluidas de E2 & etiqueta \\
    Anotación en contorno hueco       & 1 de \num{1019} regiones examinadas & anecdótico \\
    \bottomrule
  \end{tabular}
\end{table}

La subdivisión afecta al \cPctSubdiv{} de las escenas que cuentan alguna roca —aquellas en que
el número de rocas supera el de componentes conectadas— y esas escenas aportan \cRocasSubdiv{}
de las \cNRocas{} rocas, de modo que es la limitación algorítmica cuantitativamente más
importante. No toda subdivisión es un error, porque separar rocas unidas es precisamente la
función de la división de aguas; la Figura~\ref{fig:subdivision} ilustra el caso en que sí lo
es, sobre un afloramiento continuo.

\begin{figure}[H]
  \centering
  \includegraphics[width=0.9\textwidth]{Figura_22_subdivision_afloramiento.png}
  \caption{Subdivisión de un afloramiento continuo por la división de aguas.}
  \label{fig:subdivision}
\end{figure}

Los artefactos de anotación —escenas con menos del \SI{5}{\percent} de superficie etiquetada en
las que la roca mayor ocupa casi toda el área válida— son el falso positivo más llamativo. La
Figura~\ref{fig:diag-artefacto} muestra un ejemplo. Las \cNArtefacto{} escenas que cumplen ese
criterio tienen bandera \texttt{mostly\_null} y quedan fuera de la población de E2; dentro de
ella hay \cNArtefactoOk. Es la razón principal para excluir esa bandera del conteo.

\begin{figure}[H]
  \centering
  \includegraphics[width=\textwidth]{diag_artefacto_anotacion.png}
  \caption[Artefacto de anotación]{Escena con solo el \SI{2.0}{\percent} de los píxeles
    etiquetados, marcados en su totalidad como roca grande. Tiene bandera
    \texttt{mostly\_null} y no forma parte de la población de E2.}
  \label{fig:diag-artefacto}
\end{figure}

La anotación en forma de contorno hueco se midió sobre una muestra de \num{600} escenas con
roca grande y resultó \textbf{rara}: una región de las \num{1019} examinadas. Se reporta por
completitud, pero no justifica un filtro adicional.

\subsection{El límite del conteo: representación semántica y granularidad}
\label{sec:limite-conteo}

Un contraste adicional completa el diagnóstico. Los tres conjuntos de máscaras de experto que
el propio AI4Mars incluye corresponden a las \textbf{mismas} \cNGold{} imágenes con exigencias
de acuerdo crecientes entre especialistas. La proporción de escenas con roca grande disminuye
a medida que se exige mayor acuerdo (Tabla~\ref{tab:gold}).

\begin{table}[H]
  \centering
  \caption[Presencia de roca grande según la fuente de etiqueta]{Presencia de roca grande
    según la fuente de etiqueta y el nivel de acuerdo exigido. La cobertura mediana se
    calcula aquí sobre todas las escenas del conjunto —no solo sobre las que contienen
    roca—, por lo que difiere de la reportada en la Tabla~\ref{tab:validacion}.}
  \label{tab:gold}
  \begin{tabular}{lrrrr}
    \toprule
    \textbf{Fuente} & \textbf{Escenas} & \textbf{Con roca grande} & \textbf{Conteo medio}
    & \textbf{Cobertura mediana} \\
    \midrule
    Colaborativas       & \cNEscenas & \cPctBigAny  & \cColabRocas    & \cColabCobMed \\
    Experto, nivel 1    & \cNGold    & \cGoldUnoBig  & \cGoldUnoRocas  & \cGoldUnoCobMed \\
    Experto, nivel 2    & \cNGold    & \cGoldDosBig  & \cGoldDosRocas  & \cGoldDosCobMed \\
    Experto, nivel 3    & \cNGold    & \cGoldTresBig & \cGoldTresRocas & \cGoldTresCobMed \\
    \bottomrule
  \end{tabular}
\end{table}

Como las imágenes son las mismas en los tres niveles, la caída no puede deberse a diferencias de
terreno: indica que los especialistas discrepan sobre qué píxeles son roca grande, y que esos
píxeles desaparecen al exigir consenso. La frontera entre roca grande y lecho rocoso es, por
tanto, ambigua también para los expertos.

La evidencia indica que la principal limitación observada para el conteo de instancias proviene
de la representación semántica y de la granularidad de las etiquetas, más que de una falta de
resolución evidente de las imágenes. Una máscara de clases no distingue instancias; la clase
lecho rocoso absorbe la mayoría de los bloques que un observador contaría, y cuando la clase
roca grande aparece, a menudo agrupa varios bloques en un solo polígono. No es posible, con
todo, aislar la taxonomía como causa única. Intervienen también las reglas de agregación y de
consenso con que se fusionan las anotaciones, la perspectiva y la escala variable de una
cámara de navegación, la oclusión entre bloques y la propia parametrización del procedimiento.

\subsection{Síntesis: por qué contar rocas con estas máscaras no es sencillo}
\label{sec:sintesis-conteo}

El resultado negativo del conteo no procede de un único fallo, sino de varios obstáculos que se
superponen y que ningún ajuste del procedimiento resuelve por sí solo:

\begin{enumerate}
  \item \textbf{La clase que se cuenta es minoritaria.} La roca grande aporta el \cParteBig{} de
    la roca etiquetada; el resto, incluidos muchos bloques que un observador contaría, se anotó
    como lecho rocoso y queda fuera del conteo por definición.
  \item \textbf{La anotación de esa clase no es exhaustiva.} En \cVNoExhN{} de las \cVN{}
    escenas de la validación, la roca grande es menos del \SI{20}{\percent} de la roca
    etiquetada: se marcaron algunas rocas, no todas.
  \item \textbf{La geometría de la máscara no siempre separa instancias.} Cuando un polígono
    envuelve un campo de bloques contiguos, no hay estrangulamientos por donde cortar, y la
    información que separaría los bloques está en la imagen, no en la máscara.
  \item \textbf{La calibración enfrenta un dilema.} Contener la sobresegmentación de las losas
    continuas exige un criterio de semillas más estricto, y ese criterio rebajó el conteo de
    forma general (Sección~\ref{sec:calibracion}). Ninguna de las \cSenN{} combinaciones de
    parámetros ensayadas elimina el subconteo.
  \item \textbf{Usar la imagen no lo resuelve de forma demostrable.} Los relieves derivados de la
    imagen corrigen la dirección del error pero tienden a sobreestimar, y ninguno mejora el
    acuerdo con el observador una vez corregido el número de métodos probados.
\end{enumerate}

\subsection{Síntesis: inconsistencias del conjunto de datos}
\label{sec:inconsistencias}

El examen de las máscaras reveló además varias inconsistencias en el propio conjunto de datos,
que condicionan cualquier indicador derivado de él y cualquier evaluación que lo tome como
referencia. La Tabla~\ref{tab:inconsistencias} las reúne con su evidencia y su consecuencia.

\begin{table}[H]
  \centering
  \small
  \caption[Inconsistencias del conjunto de datos]{Inconsistencias del conjunto de datos AI4Mars
    (cámara de navegación de \textit{Curiosity}) identificadas en este trabajo.}
  \label{tab:inconsistencias}
  \begin{tabular}{p{3.6cm}p{5.4cm}p{5.0cm}}
    \toprule
    \textbf{Inconsistencia} & \textbf{Evidencia} & \textbf{Consecuencia} \\
    \midrule
    Composición distinta entre las escenas de entrenamiento y las de especialista
      & Medida por un mismo instrumento, cobertura mediana del \cFARegMedExp{} frente al
        \cFARegMedColab; el \cDescImgPct{} de la diferencia entre fuentes procede de las
        imágenes (Sección~\ref{sec:experto})
      & Una caída de desempeño al pasar de un conjunto al otro mezcla el cambio de anotación
        con el de terreno \\
    Criterio de anotación distinto entre voluntarios y especialistas
      & Componente de anotación de \cDescAnot{} puntos \cDescAnotIC, entre dos y cuatro según el
        instrumento
      & Las coberturas son relativas a la fuente que anotó \\
    Frontera ambigua entre roca grande y lecho rocoso
      & Sobre las mismas \cNGold{} imágenes, las escenas con roca grande pasan del
        \cGoldUnoBig{} al \cGoldTresBig{} al exigir acuerdo entre especialistas
        (Sección~\ref{sec:limite-conteo})
      & El conteo depende de qué se decidió llamar roca grande \\
    Anotación no exhaustiva de la roca grande
      & En \cVNoExhN{} de \cVN{} escenas, menos del \SI{20}{\percent} de la roca etiquetada;
        mediana del \cVNoExhMed{} (Sección~\ref{sec:anotacion-parcial})
      & El conteo no puede estimar las rocas presentes \\
    Granularidad variable del trazado
      & Polígonos que delimitan cada bloque junto a polígonos que envuelven campos enteros
        (Figura~\ref{fig:mecanismo})
      & El mismo procedimiento acierta o subcuenta según cómo se trazó \\
    Escenas casi sin etiquetar
      & \cFlagNull{} escenas con más del \SI{95}{\percent} sin etiqueta; \cNNullBig{} con roca
        grande, entre ellas los artefactos de anotación
      & Deben excluirse del conteo \\
    \bottomrule
  \end{tabular}
\end{table}

Ninguna de estas inconsistencias invalida el conjunto para el propósito con que se creó
—entrenar segmentadores de terreno para la navegación—, pero todas limitan su uso como medida
cuantitativa del terreno y como referencia de evaluación. Documentarlas es, a juicio de este
trabajo, su aporte más transferible.

\subsection{Limitaciones}

Además de las señaladas en cada apartado, conviene reunir las principales.

\paragraph{Ausencia de escala métrica.} Los tamaños son relativos al campo de visión. Sin la
geometría de la cámara o la información estereoscópica no es posible expresarlos en
centímetros ni construir distribuciones tamaño--frecuencia comparables en magnitud con la
literatura. La cámara de navegación es un par estéreo, pero el subconjunto disponible es casi
enteramente monocular.

\paragraph{Un solo observador.} La validación del conteo se realizó con un único observador, de
modo que no permite separar el desacuerdo atribuible al procedimiento de la variabilidad entre
personas. Todas las conclusiones sobre H2 son exploratorias; su confirmación exige repetir la
validación con al menos dos observadores independientes y reportar primero el acuerdo entre
ellos.

\paragraph{Muestras distintas en el contraste de fuentes.} Las máscaras colaborativas y las de
experto cubren imágenes distintas. El instrumento común permite separar ambos efectos, pero
descansa en un segmentador entrenado con etiquetas colaborativas, lo que hace aproximada la
componente atribuida a la anotación.

\paragraph{Dominio del segmentador.} El segmentador se entrenó y evaluó sobre escenas con al
menos un \SI{20}{\percent} de superficie etiquetada; fuera de ese dominio su error se triplica.
Su reparto por bloques temporales evita la fuga entre imágenes casi idénticas, pero los bloques
de prueba cubren solo una parte de la misión.

\paragraph{Subdivisión residual y subconteo.} El criterio de prominencia redujo la
fragmentación sin eliminarla: el \cPctSubdiv{} de las escenas que cuentan presenta más rocas
que componentes conectadas. La validación muestra, sin embargo, que el sesgo neto va en sentido
opuesto y es de mayor magnitud; ambos efectos coexisten en escenas distintas y no se cancelan.

\paragraph{Alcance del subconjunto.} Los resultados corresponden a una cámara de un rover en una
región. No se extienden sin más a otras misiones, otras cámaras ni otros terrenos.
